\begin{problem}{TST 2026 - Bài 2}{standard input}{standard output}{1 second}{256 megabytes}

Có một đống bóng gồm $N$ quả, các quả bóng được đánh số từ $0$ đến $N-1$. Mỗi quả bóng trong đống có thể có màu từ $0$ đến $K-1$, trong đó quả bóng thứ $i$ ($0\le i\le N-1$) có màu $A_i$ ($0\le A_i\le K-1$). Đảm bảo với mọi $j$ từ $0$ đến $K-1$, tồn tại \textbf{ít nhất} một quả bóng trong đống bóng có màu $j$.

Bạn được biết trước giá trị của $N$, nhưng bạn lại \textbf{không biết chính xác} giá trị của $K$. Bạn không biết trạng thái màu của các quả bóng trông thế nào, và bạn muốn tìm cách xác định nó. Để làm được việc này, bạn được phép thực hiện một số câu hỏi. Với mỗi câu hỏi:
\begin{itemize}
   \item Bạn chọn một tập con gồm một số quả bóng trong $N$ quả.
   \item Sau đó, bạn sẽ được biết \textbf{số lượng màu khác nhau} xuất hiện trong các quả bóng được chọn có lớn hơn hoặc bằng $\lfloor \frac{K+1}2\rfloor$ hay không (trong đó, $\lfloor x\rfloor$ là số nguyên gần nhất nhỏ hơn hoặc bằng $x$).
\end{itemize}

\textbf{Yêu cầu:} Hãy viết một chương trình tương tác mô phỏng quá trình xác định trạng thái màu của các quả bóng. Vì có thể có nhiều cấu hình cho ra cùng một kết quả với mọi câu hỏi so với cấu hình gốc, bạn chỉ cần xác định một cấu hình \textbf{đẳng cấu} với cấu hình gốc. Cụ thể:
\begin{itemize}
   \item Gọi $B$ là cấu hình (trạng thái màu của các quả bóng) mà bạn dự đoán. Khi đó với mọi cặp số $(i,j)$ ($0\le i,j\le N-1,i\ne j$):
   \begin{itemize}
      \item Nếu $A_i=A_j$ thì $B_i=B_j$.
      \item Nếu $A_i\ne A_j$ thì $B_i\ne B_j$.
   \end{itemize}
   \item \textbf{Ví dụ:} Hai cấu hình \texttt{{0, 1, 2, 0, 0, 1}} và \texttt{{1, 2, 0, 1, 1, 2}} được xem như là \textbf{đẳng cấu} với nhau.
\end{itemize}

\InputFile
\textbf{Chi tiết cài đặt}

Thí sinh cần cài đặt hàm sau:

\begin{lstlisting}
vector<int> solve(int N);
\end{lstlisting}

\begin{itemize}
   \item $N$: Số quả bóng trong đống bóng cho trước.
   \item Hàm này cần trả về một mảng $B$ gồm $N$ phần tử đại diện cho cấu hình (trạng thái màu của các quả bóng) mà bạn dự đoán. Cấu hình của bạn cần phải đảm bảo các điều kiện như sau:
   \begin{itemize}
      \item $0\le B_i\le K-1$ với mọi $i$ từ $0$ đến $N-1$.
      \item $B$ \textbf{đẳng cấu} với $A$ ($A$ là mảng gồm $N$ phần tử đại diện cho trạng thái màu gốc của $N$ quả bóng).
   \end{itemize}
   \item \textbf{Lưu ý:} Bạn không được cho trước giá trị chính xác của $K$.
   \item Hàm này được gọi nhiều lần, mỗi lần gọi tương ứng với một trường hợp thử nghiệm.
\end{itemize}

Trong hàm \texttt{solve()}, thí sinh được phép gọi hàm sau:

\begin{lstlisting}
bool ask(vector<int> pos);
\end{lstlisting}

\begin{itemize}
   \item $pos$: Một mảng bao gồm chỉ số của tất cả quả bóng được chọn. Cần phải đảm bảo các phần tử trong mảng $pos$ là các chỉ số quả bóng hợp lệ, và các giá trị trong mảng $pos$ phải đôi một phân biệt.
   \item Hàm này sẽ trả về \texttt{true} nếu như số màu khác nhau trong đống bóng được chọn lớn hơn hoặc bằng $\lfloor \frac{K+1}2\rfloor$, ngược lại hàm sẽ trả về \texttt{false}.
   \item Với mỗi trường hợp thử nghiệm, hàm \texttt{solve()} được gọi hàm \texttt{ask()} tối đa $14000$ lần.
\end{itemize}

\textbf{Lưu ý:}
\begin{itemize}
   \item Chương trình của bạn cần phải khai báo thư viện \texttt{"colorfullib.h"} ở đầu.
   \item Trong chương trình, thí sinh được phép cài đặt các biến, các hàm toàn cục, nhưng không được phép có hàm \texttt{main()}.
   \item Chương trình của bạn không được tương tác trực tiếp với đầu vào chuẩn (stdin) và đầu ra chuẩn (stdout).
\end{itemize}

\textbf{Ràng buộc}

\begin{itemize}
   \item $3\le K\le 32$
   \item $K\le N\le 256$
   \item $\sum_N\le 2\times 10^4$
\end{itemize}

\Scoring
\begin{center}
  \begin{tabular}{|c|c|l|} \hline
    \textbf{Subtask} & \textbf{Điểm} & \textbf{Giới hạn} \\ \hline
    1 & $10$ & $N\le 10;K\le 3$ \\ \hline
    2 & $20$ & $N\le 10$ \\ \hline
    3 & $70$ & Không có ràng buộc gì thêm \\ \hline
  \end{tabular}
\end{center}

\textbf{Cách tính điểm}

Nếu trong bất kỳ trường hợp thử nghiệm nào, chương trình của bạn trả về cấu hình không thỏa mãn yêu cầu đề bài thì bạn sẽ không nhận được điểm cho test đó.

Ngược lại, gọi $P$ là số lần gọi hàm \texttt{ask()} tối đa trong tất cả lần gọi hàm \texttt{solve()} của một test. Khi đó điểm của bạn được tính như sau:

\begin{center}
  \begin{tabular}{|c|c|} \hline
    \textbf{Giá trị $P$} & \textbf{Phần trăm số điểm} \\ \hline
    $P \le 1400$ & $100\%$ \\ \hline
    $1400 \lt P \le 14000$ & $\frac{1400}P\times 100\%$ \\ \hline
    $P \gt 14000$ & $0\%$ \\ \hline
  \end{tabular}
\end{center}

Điểm của một subtask bằng phần trăm điểm tối thiểu của một test trong subtask đó nhân với số điểm tối đa của subtask.


\end{problem}

